% ==============================================================================
% ch03.tex -- Chapter 3: The Four Basic Moduli Problems for Elliptic Curves:
%             Sorites
%   Source: book pp. 98--106 (physical pp. 109--117)
% ==============================================================================
\chapter{The four basic moduli problems for elliptic curves: Sorites}
\label{ch:3}
% The printed running head omits the subtitle (original: "3. THE FOUR BASIC
% MODULI PROBLEMS FOR ELLIPTIC CURVES"), which would otherwise collide with
% the folio.  The table of contents keeps the full title.
\markboth{3.\ \MakeUppercase{The four basic moduli problems for elliptic curves}}{}

Let $S$ be a scheme, $E/S$ an elliptic curve over $S$, and $N \ge 1$ an integer. We will make use of the Drinfeldian ideas of the first chapter to define four sorts of ``level structure'' on $E/S$, which correspond to the classical notions of $\Gamma(N)$, $\Gamma_1(N)$, ``balanced''-$\Gamma_1(N)$, and $\Gamma_0(N)$ structures (and which coincide with these notions when $N$ is invertible on $S$).

% ------------------------------------------------------------------------------
% Section 3.1
% ------------------------------------------------------------------------------
\section{$\Gamma(N)$-structures}
\label{sec:3.1}

A $\Gamma(N)$-structure on $E/S$ (also called a ``full level $N$ structure'', or a ``Drinfeld basis of $E[N]$'') is a group homomorphism
\begin{equation}\label{eq:3.1.1}
\phi : (\Z/N\Z)^2 \longrightarrow E[N](S)
\end{equation}
which is a ``generator'' of $E[N]$ in the sense of Chapter~\ref{ch:1}, \ref{sec:1.5}. Explicitly, this means that we have an equality of effective Cartier divisors in $E$ :
\begin{equation}\label{eq:3.1.2}
E[N] = \sum_{a, b \bmod N} [\phi(a,b)]
\end{equation}
or equivalently that the $N^2$ sections $\phi(a,b)$ of $E[N](S)$ form a ``full set of sections.'' The points
\[
P = \phi(1,0), \qquad Q = \phi(0,1).
\]
in $E[N](S)$ are the corresponding ``Drinfeld basis'' of $E[N]$.

% ------------------------------------------------------------------------------
% Section 3.2
% ------------------------------------------------------------------------------
\section{$\Gamma_1(N)$-structures}
\label{sec:3.2}

A $\Gamma_1(N)$-structure on $E/S$, also called a point of ``exact order $N$'' in $E(S)$, is a homomorphism
\begin{equation}\label{eq:3.2.1}
\phi : \Z/N\Z \longrightarrow E[N](S)
\end{equation}
which is a ``$\Z/N\Z$-structure on $E[N]/S$'' in the sense of \ref{item:1.5.1}. Explicitly, this means that the effective Cartier divisor in $E$
\begin{equation}\label{eq:3.2.2}
\sum_{a \bmod N} [\phi(a)]
\end{equation}
is a subgroup-scheme of $E$. The point
\[
P = \phi(1)
\]
in $E[N](S)$ is the corresponding point of ``exact order $N$.''

Equivalently, a $\Gamma_1(N)$-structure on $E/S$ is an $N$-isogeny of elliptic curves over $S$
\begin{equation}\label{eq:3.2.3}
E \xrightarrow{\ \pi\ } E'
\end{equation}
(by definition, a homomorphism which is finite locally free of degree $N$) together with a generator of $\Ker \pi$, i.e., a point
\begin{equation}\label{eq:3.2.4}
P \in (\Ker \pi)(S) \subset E[N](S)
\end{equation}
such that the corresponding homomorphism
\[
\begin{aligned}
\phi : \Z/N\Z &\longrightarrow \Ker \pi \\
a &\longmapsto aP
\end{aligned}
\]
is a ``generator'' of $\Ker \pi$ in the sense of the \ref{item:1.4.1}, i.e., such that we have an equality of effective Cartier divisors in $E$
\begin{equation}\label{eq:3.2.5}
\Ker \pi = \sum_{a \bmod N} [aP].
\end{equation}

% ------------------------------------------------------------------------------
% Section 3.3
% ------------------------------------------------------------------------------
\section{Balanced $\Gamma_1(N)$-structures}
\label{sec:3.3}

A balanced $\Gamma_1(N)$-structure on $E/S$ is a diagram
\begin{equation}\label{eq:3.3.1}
P;\quad E \xleftrightarrow[\ \pi^t\ ]{\ \pi\ } E';\quad P'
\end{equation}
in which $E'$ is an elliptic curve over $S$, $\pi$ is an $N$-isogeny, $\pi^t$ is the dual $N$-isogeny ($\pi^t \pi = N_E$, $\pi \pi^t = N_{E'}$), $P \in (\Ker \pi)(S)$ is a generator of $\Ker \pi$, and $P' \in (\Ker \pi^t)(S)$ is a generator of $\Ker \pi^t$.

Less symmetrically, we can (by \cite[III, Exp.~V, 4.1]{SGA}) describe a balanced $\Gamma_1(N)$-structure on $E/S$ purely ``on $E$'' as consisting of an f.p.p.f. short exact sequence of group-schemes on $S$
\begin{equation}\label{eq:3.3.2}
0 \longrightarrow K \longrightarrow E[N] \longrightarrow K' \longrightarrow 0,
\end{equation}
with $K$ and $K'$ both locally free of rank $N$, together with points
\begin{equation}\label{eq:3.3.3}
P \in K(S), \qquad P' \in K'(S)
\end{equation}
which are generators of $K$ and $K'$ respectively.

% ------------------------------------------------------------------------------
% Section 3.4
% ------------------------------------------------------------------------------
\section{$\Gamma_0(N)$-structures}
\label{sec:3.4}

A $\Gamma_0(N)$-structure on $E/S$ is an $N$-isogeny
\begin{equation}\label{eq:3.4.1}
E \xrightarrow{\ \pi\ } E'
\end{equation}
which is cyclic in the sense that locally f.p.p.f. on $S$, the kernel $\Ker(\pi)$ admits a generator. Equivalently, a $\Gamma_0(N)$-structure on $E/S$ is a finite flat subgroup-scheme
\begin{equation}\label{eq:3.4.2}
K \subset E[N],
\end{equation}
locally free of rank $N$, which is cyclic in the sense that locally f.p.p.f. on $S$, it admits a generator.

% ------------------------------------------------------------------------------
% Section 3.5
% ------------------------------------------------------------------------------
\section{Factorization into prime powers}
\label{sec:3.5}

\begin{plem}\label{lem:3.5.1}
Suppose that $N = AB$ with $A$ and $B$ relatively prime. Then for any elliptic curve $E/S$, we have functorial isomorphisms
\[
\small\begin{array}{@{}l@{}}
\displaystyle \Gamma(N)\text{-Str}(E/S) \xrightarrow{\ \sim\ } \Gamma(A)\text{-Str}(E/S) \times \Gamma(B)\text{-Str}(E/S) \\[2pt]
\displaystyle \Gamma_1(N)\text{-Str}(E/S) \xrightarrow{\ \sim\ } \Gamma_1(A)\text{-Str}(E/S) \times \Gamma_1(B)\text{-Str}(E/S) \\[2pt]
\displaystyle \mathrm{Bal}\ \Gamma_1(N)\text{-Str}(E/S) \xrightarrow{\ \sim\ } \mathrm{Bal}\ \Gamma_1(A)\text{-Str}(E/S) \times \mathrm{Bal}\ \Gamma_1(B)\text{-Str}(E/S) \\[2pt]
\displaystyle \Gamma_0(N)\text{-Str}(E/S) \xrightarrow{\ \sim\ } \Gamma_0(A)\text{-Str}(E/S) \times \Gamma_0(B)\text{-Str}(E/S).
\end{array}
\]
\end{plem}

\begin{proof}[Proof-construction]
For any abelian group $G$ killed by $N$, we have two distinct canonical isomorphisms
\begin{equation}\label{eq:3.5.1.1}
G \longrightarrow G[A] \times G[B]
\end{equation}
of $G$ with the products of its $A$-torsion and $B$-torsion subgroups. The first is the inverse of the ``sum'' map
\begin{equation}\label{eq:3.5.1.2}
G[A] \times G[B] \overset{\text{sum}}{\xrightarrow{\ \sim\ }} G,
\end{equation}
the second is
\begin{equation}\label{eq:3.5.1.3}
\begin{gathered}
G \xrightarrow{\ \sim\ } G[A] \times G[B] \\
g \longmapsto (Bg, Ag).
\end{gathered}
\end{equation}
[The second is $(A+B)$ times the first.] Applying \emph{either} of these decompositions to our structures defines isomorphisms of the kind desired as follows easily from the general theorems of factorization into prime powers of ``$A$-structures'' and ``$A$-generators'', cf., \ref{cor:1.7.3}. For definiteness, let us choose the second. Then explicitly the isomorphisms are given by
\[
\small\left\{
\begin{array}{@{}l@{\;}r@{\;}r@{}}
\text{for } \Gamma(N) : & \text{Drinfeld basis } (P,Q) \mapsto & \text{Drinfeld bases } (BP, BQ), \\
 & & \text{and } (AP, AQ) \\[5pt]
\text{for } \Gamma_1(N) : & \text{point } \underline{P} \text{ of exact order } N \mapsto & \text{points } BP, AP \text{ of} \\
 & & \text{exact orders } A, B \\[5pt]
\text{for } \Gamma_0(N) : & \text{cyclic subgroup-scheme } K \mapsto & \text{cyclic subgroups } K[A], \\
 & \text{of order } N & K[B] \text{ of orders } A, B.
\end{array}
\right.
\]

To make explicit the isomorphism for balanced $\Gamma_1(N)$-structures, we view such a structure ``asymmetrically'', as an f.p.p.f. short exact sequence together with generators of end terms
\[
\begin{tikzcd}[column sep=large, row sep=tiny]
0 \arrow[r] & K \arrow[r] & E[N] \arrow[r] & K' \arrow[r] & 0. \\
& \underset{P}{\sim} & & \underset{P'}{\sim}
\end{tikzcd}
\]
The associated balanced $\Gamma_1(A)$-structure is then
\[
\begin{tikzcd}[column sep=large, row sep=tiny]
0 \arrow[r] & K[A] \arrow[r] & E[A] \arrow[r] & K'[A] \arrow[r] & 0 \\
& \underset{BP}{\sim} & & \underset{BP'}{\sim}
\end{tikzcd}
\]
and similarly for the $\Gamma_1(B)$-structure (just interchange $A$ and $B$).
\end{proof}

% ------------------------------------------------------------------------------
% Section 3.6
% ------------------------------------------------------------------------------
\section{Relative representability}
\label{sec:3.6}

\begin{prelrep}{3.6.0}\label{thm:3.6.0}
Fix an elliptic curve $E/S$, and an integer $N \ge 1$. Consider the three functors on $(\Sch/S)$ defined by
\[
T \mapsto \begin{cases}
\Gamma(N)\text{-structures on } E_T/T \\
\Gamma_1(N)\text{-structures on } E_T/T \\
\text{balanced } \Gamma_1(N)\text{-structures on } E_T/T.
\end{cases}
\]
Each of these functors is represented by a finite $S$-scheme.
\end{prelrep}

\begin{proof}
For the first two functors, this is a special case of the representability results of the previous chapter, in whose notations the first two functors are respectively
\[
(\Z/N\Z)^2\text{-Gen}(E[N]/S) \quad \text{and} \quad \Z/N\Z\text{-Str}(E/S).
\]
The third functor is represented by a scheme finite over $\Z/N\Z\text{-Str}(E/S)$ namely the scheme
\[
\Z/N\Z\text{-Gen}\Bigl(K'\Big/\bigl(\Z/N\Z\text{-Str}(E/S)\bigr)\Bigr),
\]
where $K'$ is the tautological quotient of $E[N]$ by the subgroup $K$ specified by a $\Z/N\Z$-structure on $E[N]$.
\end{proof}

\begin{prem}{3.6.1}\label{rem:3.6.1}
We will see later that the functor of $\Gamma_0(N)$-structures on $E/S$ is also representable by a finite $S$-scheme, but this fact lies much deeper.
\end{prem}

If we apply the factorization lemma \ref{cor:1.7.3} to the finite schemes representing these functors, we obtain

\begin{pcor}{3.6.2}\label{cor:3.6.2}
If $N = AB$ with $(A,B) = 1$, then the constructions \eqref{eq:3.5.1.3} define canonical isomorphisms of finite $S$-schemes
\[
\Gamma(N)\text{-Str}(E/S) \xrightarrow{\ \sim\ } \Gamma(A)\text{-Str}(E/S) \underset{S}{\times} \Gamma(B)\text{-Str}(E/S)
\]
and similarly for $\Gamma_1$ and for balanced $\Gamma_1$.
\end{pcor}

\begin{prem}{3.6.2.1}\label{rem:3.6.2.1}
Similarly, the functor $\Gamma_0(N)$-Str$(E/S)$ on $(\Sch/S)$ defined by
\[
T \mapsto \Gamma_0(N)\text{-structures on } E_T/T
\]
admits a canonical product decomposition, as functor on $(\Sch/S)$,
\[
\Gamma_0(N)\text{-Str}(E/S) \xrightarrow{\ \sim\ } \Gamma_0(A)\text{-Str}(E/S) \underset{S}{\times} \Gamma_0(B)\text{-Str}(E/S).
\]
This decomposition reduces the problem of proving representability for $\Gamma_0(N)$ to the case when $N$ is a power of a prime.
\end{prem}

\begin{pprop}{3.6.3}\label{prop:3.6.3}
Let $E$ and $E'$ be elliptic curves over a base scheme $S$. Suppose that we are given an isomorphism of $S$-group-schemes
\[
E[N] \xrightarrow{\ \sim\ } E'[N].
\]
Then we have induced isomorphisms of functors on $(\Sch/S)$
\[
\Gamma(N)\text{-Str}(E'/S) \xrightarrow{\ \sim\ } \Gamma(N)\text{-Str}(E/S)
\]
and similarly for $\Gamma_1(N)$, balanced $\Gamma_1(N)$, and $\Gamma_0(N)$.
\end{pprop}

\begin{proof}
Indeed all of these sorts of structures depend only on the underlying $S$-group-scheme $E[N]$.
\end{proof}

% ------------------------------------------------------------------------------
% Section 3.7
% ------------------------------------------------------------------------------
\section{The situation when $N$ is invertible}
\label{sec:3.7}

\begin{theorem}\label{thm:3.7.1}
Let $N \ge 1$ be an integer, $S$ a scheme on which $N$ is invertible (i.e., $S$ is a $\Z[1/N]$-scheme), and $E/S$ an elliptic curve. Consider the four functors on $(\Sch/S)$ given by
\[
T \mapsto \begin{cases}
\Gamma(N)\text{-structures on } E_T/T \\
\Gamma_1(N)\text{-structures on } E_T/T \\
\text{balanced } \Gamma_1(N)\text{-structures on } E_T/T \\
\Gamma_0(N)\text{-structures on } E_T/T.
\end{cases}
\]
Each is represented by a finite etale $S$-scheme.
\end{theorem}

\begin{proof}
Because $N$ is invertible on $S$, the group-scheme $E[N]$ is finite etale over $S$, locally (etale) isomorphic to $(\Z/N\Z)^2$ (cf.\ \ref{thm:2.3.1}). The assertion for $\Gamma(N)$ (resp.\ for $\Gamma_1(N)$) therefore results immediately from \ref{prop:1.6.7} applied to $E[N]/S$ (resp.\ from \ref{prop:1.6.4} applied to $A = \Z/N\Z$, $C = E$). The assertion for balanced $\Gamma_1(N)$ results by combining \ref{prop:1.6.4} for $A = \Z/N$, $C = E$ over $S$ with \ref{prop:1.6.7} applied to the base scheme $A$-Str$(E/S)$, the group $A = \Z/N\Z$, and the quotient group-scheme $G = E[N] \bmod K_{\mathrm{univ}}$ inside the curve $E \bmod K_{\mathrm{univ}}$, where $K_{\mathrm{univ}}$ is the tautological subgroup of $E[N]$ specified by the universal $\Z/N\Z$-structure.

For the $\Gamma_0(N)$-problem, we argue as follows. The functor in question is an etale sheaf (etale descent \cite[I, 7.9]{SGA} for finite locally free subgroup-schemes of $E[N]$, and the fact that the notion of cyclicity is by definition local for the f.p.p.f. topology, so a fortiori for the etale topology). Therefore by etale descent it suffices to show that after a surjective etale base change $S' \to S$, the $\Gamma_0(N)$ problem becomes representable by a finite etale $S'$-scheme. By \ref{thm:2.3.1}, we may make such a base change and reduce to the case when $E[N]$ is the constant group-scheme $(\Z/N\Z)^2$. In this case, we claim that the $\Gamma_0(N)$ functor is represented by the constant $S$-scheme
\[
S \times \{\text{the set of cyclic subgroups of order } N \text{ in } (\Z/N\Z)^2\}.
\]
To show this, it suffices to check that over any connected scheme $T$, any finite locally-free closed subgroup-scheme $K$ of rank $N$ of a constant (e.g., $(\Z/N\Z)^2$) finite locally-free group-scheme $G$ over $T$ is itself constant. For once we know this, it follows from \ref{prop:1.10.12} (3) that our constant group-scheme $K$ is cyclic in the sense of f.p.p.f.-locally admitting a generator if and only if the abstract group $K(T)$ is a cyclic group of order $N$.

To prove that $K$ is constant, we may reduce easily to the case where $T$ is connected and noetherian (\cite[IV, 8.9.1, 11.2.6.1]{EGA}). Then $K$ is finite etale over $T$, because it is both ``net'' \cite[I, 4.3.1--2]{DeGa}, being a closed subscheme of an etale (because constant) $T$-scheme, and finite flat over $T$.

Finally, over any connected, locally noetherian scheme $T$, any finite etale $T$-subscheme $K$ of a constant finite etale $T$-scheme $G$ is itself constant. To see this, pick a geometric point $t$ of $T$. In view of the standard theory of the fundamental group (\cite[I, Exp.~V, \S7]{SGA}), a finite etale $T$-scheme $K$ is constant if and only if the action of $\pi_1(T,t)$ on the fiber $K(t)$ is trivial, and by assumption, $K(t)$ is a $\pi_1(T,t)$-stable subset of the finite $\pi_1(T,t)$-set $G(t)$, which has trivial action of $\pi_1$.
\end{proof}

\begin{corollary}\label{cor:3.7.2}
Let $N \ge 1$ be an integer, $S$ a scheme on which $N$ is invertible, and $E/S$ an elliptic curve. Suppose that $E[N]$ is $S$-isomorphic to the constant group-scheme $(\Z/N\Z)^2$ (a condition which by \ref{thm:3.7.1} is always fulfilled after a finite etale surjective base-change $S' \to S$). Then the four functors on $(\Sch/S)$ defined by
\[
T \mapsto \left\{\begin{array}{@{}l@{}}
\Gamma(N)\text{-structures} \\
\Gamma_1(N)\text{-structures} \\
\text{balanced } \Gamma_1(N)\text{-structures} \\
\Gamma_0(N)\text{-structures}
\end{array}\right\} \text{ on } E_T/T
\]
are represented by the constant $S$-schemes
\[
S \times \{\text{the set of } \Z/N\Z\text{-bases of } (\Z/N\Z)^2\}
\]
\[
S \times \{\text{the set of elements } P \text{ of } (\Z/N\Z)^2 \text{ having exact order } N\}
\]
\[
\small S \times \left\{\begin{array}{@{}l@{}}
\text{the set of triples } (K, P, P') \text{ consisting of a cyclic subgroup } K \\
\text{of } (\Z/N\Z)^2 \text{ of order } N, \text{ a generator } P \text{ of } K, \text{ and a generator} \\
P' \text{ of the quotient } (\Z/N\Z)^2 \text{ modulo } K
\end{array}\right\}
\]
\[
S \times \{\text{the set of cyclic subgroups of } (\Z/N\Z)^2 \text{ of order } N\}.
\]
\end{corollary}

\begin{proof}
For $\Gamma(N)$, this is obvious from \ref{prop:1.10.12} (3). For $\Gamma_0(N)$, this was proven in the course of proving \ref{thm:3.7.1}. From the result for $\Gamma_0(N)$, the results for $\Gamma_1(N)$ (resp.\ for balanced $\Gamma_1(N)$) follow immediately from \ref{prop:1.10.12} (3) applied to $A = \Z/N\Z$ and to $G = K$ (resp.\ to $G = K$ and to $G = (\Z/N\Z)^2$ modulo $K$).
\end{proof}
